% IU81001C
% Multiresistente Keime
% 14.0
% 2013-11-30
%
:ref:`m-mrsa-transport`%
\label{proc:IU81001C}
\begin{itemize}
  -   Informationen über MRSA-Trägerschaft sowie Art/Ort der MRSA-Besiedelung
  -   Patientenvorbereitung und Transport
  \begin{itemize}
    -   Wunden müssen verbunden und gut abgedeckt sein
    -   Bei Besiedelung der Atemwege trägt der Patient einen Mund-/ Nasenschutz
    -   Vor dem Transport führt der Patient eine hygienische
    Händedesinfektion durch
  \end{itemize}
  -   Allgemeine Hygienemaßnahmen
  \begin{itemize}
    -   Einweghandschuhe und Schutzkittel
    -   Nach Ablegen sofortige hygienische Händedesinfektion
    \begin{itemize}
      -   Nach Abschluss des Patiententransportes sind alle Materialien,
      Geräte, Instrumente und Flächen, welche direkten Kontakt mit dem
      Patienten hatten, gemäß Hygieneplan zu
      desinfizieren.
      -   Alle waagerechten Oberflächen des Fahrzeuginnenraumes sind mit
      einem Flächendesinfektionsmittel in Form einer
      Scheuerwischdesinfektion desinfizierend zu reinigen.
      -   Abschließend führt das Personal eine hygienische
      Händedesinfektion durch.
      -   An Keimvehikel denken (Kugelschreiber, Haltegriffe, Stethoskop \ldots )!
    \end{itemize}
  \end{itemize}
  -   Nach Abschluss der vorangegangenen Hygiene- und
  Desinfektionsmaßnahmen ist das Fahrzeug wieder einsatzbereit.
\end{itemize}

Für die anderen Problemkeime \Z{(VRSA, VRE, ESBL,
{\dots})} gelten sinngemäß die gleichen
Verfahrensweisen.
